\chapter{Produce a stream another decoder can read}
\MARSChapterHeading{Produce a stream another decoder can read}
\label{revision:producing-video-streams}

The previous chapter ended with selected coefficient levels and a reconstructed
picture. Neither is yet an H.263 file. The receiver needs an ordered description
of picture type, prediction, coefficient positions and signed values. A wrong
presence flag can change how every later bit is interpreted even when the
underlying transform was accurate.

This chapter completes the restricted encoder in \texttt{h263-encoder.asm}.
We will follow a field across byte boundaries, describe a block's nonzeros,
combine six blocks into a macroblock, and publish complete pictures through
ordinary file I/O. An independent parser and decoder then check what the target
program actually wrote. The final independent exercise reads bounded fields
back from a byte sequence, giving us a small foundation for an inspector.

Our profile uses standard QCIF, one initial I-picture followed by P-pictures,
one quantizer per picture, whole-pixel luma search and explicit group headers
between macroblock rows. Extended picture types, unrestricted vectors, arithmetic
coding, advanced prediction and PB pictures are disabled. These are documented
restrictions on an actual H.263 stream. A general decoder can read the output;
our teaching parser deliberately admits a narrower set than every legal H.263
stream.

\section{Keep field width separate from machine byte order}

Value five in three bits is \texttt{101}; in five bits it is \texttt{00101}.
The width supplies leading zeros that a stored integer does not remember.
Emit the most significant declared bit first. In the output array, the first
bit occupies bit seven of byte zero, followed by bit six. A new field continues
at the next free bit, including inside a partly used byte.

For a committed count $p$, the next byte index is $\lfloor p/8\rfloor$ and its
bit position is $7-(p\bmod8)$. This transmission order is independent of native
word byte order. The writer must produce identical serialized bytes on either
simulator target. Storing a native word is not a substitute for appending its
chosen field bits.

Start with bytes \texttt{A5 A5}, but only their first three bits committed.
Those valid bits are \texttt{101}. Append the nine-bit value 419,
\texttt{110100011}. The bytes become \texttt{BA 35}, with twelve valid bits.
The low four bits of the second byte still contain the old \texttt{0101}; they
are outside the append. Append four explicit zeros to reach sixteen committed
bits and the bytes become \texttt{BA 30}. Alignment changes the valid bit count;
it is not a license to overwrite an unfinished byte without recording the bits.

\begin{figure}[!ht]
  \centering
  \begin{tikzpicture}[every node/.style={font=\small},alt={Three sixteen-bit rows show A5 A5 with only prefix 101 valid, BA 35 with twelve bits valid after appending 110100011, and BA 30 with all sixteen bits valid after explicit zero alignment. Unused cells have dashed outlines and pale fill.}]
    \node[anchor=east] at (-.25,0.23) {3 valid};
    \draw[black,line width=.4pt] (0.000,-0.000) rectangle (0.420,0.460);\node at (0.210,0.230) {1};
    \draw[black,line width=.4pt] (0.420,-0.000) rectangle (0.840,0.460);\node at (0.630,0.230) {0};
    \draw[black,line width=.4pt] (0.840,-0.000) rectangle (1.260,0.460);\node at (1.050,0.230) {1};
    \draw[black,dashed,fill=black!10,line width=.4pt] (1.260,-0.000) rectangle (1.680,0.460);\node at (1.470,0.230) {0};
    \draw[black,dashed,fill=black!10,line width=.4pt] (1.680,-0.000) rectangle (2.100,0.460);\node at (1.890,0.230) {0};
    \draw[black,dashed,fill=black!10,line width=.4pt] (2.100,-0.000) rectangle (2.520,0.460);\node at (2.310,0.230) {1};
    \draw[black,dashed,fill=black!10,line width=.4pt] (2.520,-0.000) rectangle (2.940,0.460);\node at (2.730,0.230) {0};
    \draw[black,dashed,fill=black!10,line width=.4pt] (2.940,-0.000) rectangle (3.360,0.460);\node at (3.150,0.230) {1};
    \draw[black,dashed,fill=black!10,line width=.4pt] (3.360,-0.000) rectangle (3.780,0.460);\node at (3.570,0.230) {1};
    \draw[black,dashed,fill=black!10,line width=.4pt] (3.780,-0.000) rectangle (4.200,0.460);\node at (3.990,0.230) {0};
    \draw[black,dashed,fill=black!10,line width=.4pt] (4.200,-0.000) rectangle (4.620,0.460);\node at (4.410,0.230) {1};
    \draw[black,dashed,fill=black!10,line width=.4pt] (4.620,-0.000) rectangle (5.040,0.460);\node at (4.830,0.230) {0};
    \draw[black,dashed,fill=black!10,line width=.4pt] (5.040,-0.000) rectangle (5.460,0.460);\node at (5.250,0.230) {0};
    \draw[black,dashed,fill=black!10,line width=.4pt] (5.460,-0.000) rectangle (5.880,0.460);\node at (5.670,0.230) {1};
    \draw[black,dashed,fill=black!10,line width=.4pt] (5.880,-0.000) rectangle (6.300,0.460);\node at (6.090,0.230) {0};
    \draw[black,dashed,fill=black!10,line width=.4pt] (6.300,-0.000) rectangle (6.720,0.460);\node at (6.510,0.230) {1};
    \draw[line width=.9pt] (3.36,-0.000)--(3.36,0.460);
    \node[anchor=east] at (-.25,-0.62) {12 valid};
    \draw[black,line width=.4pt] (0.000,-0.850) rectangle (0.420,-0.390);\node at (0.210,-0.620) {1};
    \draw[black,line width=.4pt] (0.420,-0.850) rectangle (0.840,-0.390);\node at (0.630,-0.620) {0};
    \draw[black,line width=.4pt] (0.840,-0.850) rectangle (1.260,-0.390);\node at (1.050,-0.620) {1};
    \draw[black,line width=.4pt] (1.260,-0.850) rectangle (1.680,-0.390);\node at (1.470,-0.620) {1};
    \draw[black,line width=.4pt] (1.680,-0.850) rectangle (2.100,-0.390);\node at (1.890,-0.620) {1};
    \draw[black,line width=.4pt] (2.100,-0.850) rectangle (2.520,-0.390);\node at (2.310,-0.620) {0};
    \draw[black,line width=.4pt] (2.520,-0.850) rectangle (2.940,-0.390);\node at (2.730,-0.620) {1};
    \draw[black,line width=.4pt] (2.940,-0.850) rectangle (3.360,-0.390);\node at (3.150,-0.620) {0};
    \draw[black,line width=.4pt] (3.360,-0.850) rectangle (3.780,-0.390);\node at (3.570,-0.620) {0};
    \draw[black,line width=.4pt] (3.780,-0.850) rectangle (4.200,-0.390);\node at (3.990,-0.620) {0};
    \draw[black,line width=.4pt] (4.200,-0.850) rectangle (4.620,-0.390);\node at (4.410,-0.620) {1};
    \draw[black,line width=.4pt] (4.620,-0.850) rectangle (5.040,-0.390);\node at (4.830,-0.620) {1};
    \draw[black,dashed,fill=black!10,line width=.4pt] (5.040,-0.850) rectangle (5.460,-0.390);\node at (5.250,-0.620) {0};
    \draw[black,dashed,fill=black!10,line width=.4pt] (5.460,-0.850) rectangle (5.880,-0.390);\node at (5.670,-0.620) {1};
    \draw[black,dashed,fill=black!10,line width=.4pt] (5.880,-0.850) rectangle (6.300,-0.390);\node at (6.090,-0.620) {0};
    \draw[black,dashed,fill=black!10,line width=.4pt] (6.300,-0.850) rectangle (6.720,-0.390);\node at (6.510,-0.620) {1};
    \draw[line width=.9pt] (3.36,-0.850)--(3.36,-0.390);
    \node[anchor=east] at (-.25,-1.47) {16 valid};
    \draw[black,line width=.4pt] (0.000,-1.700) rectangle (0.420,-1.240);\node at (0.210,-1.470) {1};
    \draw[black,line width=.4pt] (0.420,-1.700) rectangle (0.840,-1.240);\node at (0.630,-1.470) {0};
    \draw[black,line width=.4pt] (0.840,-1.700) rectangle (1.260,-1.240);\node at (1.050,-1.470) {1};
    \draw[black,line width=.4pt] (1.260,-1.700) rectangle (1.680,-1.240);\node at (1.470,-1.470) {1};
    \draw[black,line width=.4pt] (1.680,-1.700) rectangle (2.100,-1.240);\node at (1.890,-1.470) {1};
    \draw[black,line width=.4pt] (2.100,-1.700) rectangle (2.520,-1.240);\node at (2.310,-1.470) {0};
    \draw[black,line width=.4pt] (2.520,-1.700) rectangle (2.940,-1.240);\node at (2.730,-1.470) {1};
    \draw[black,line width=.4pt] (2.940,-1.700) rectangle (3.360,-1.240);\node at (3.150,-1.470) {0};
    \draw[black,line width=.4pt] (3.360,-1.700) rectangle (3.780,-1.240);\node at (3.570,-1.470) {0};
    \draw[black,line width=.4pt] (3.780,-1.700) rectangle (4.200,-1.240);\node at (3.990,-1.470) {0};
    \draw[black,line width=.4pt] (4.200,-1.700) rectangle (4.620,-1.240);\node at (4.410,-1.470) {1};
    \draw[black,line width=.4pt] (4.620,-1.700) rectangle (5.040,-1.240);\node at (4.830,-1.470) {1};
    \draw[black,line width=.4pt] (5.040,-1.700) rectangle (5.460,-1.240);\node at (5.250,-1.470) {0};
    \draw[black,line width=.4pt] (5.460,-1.700) rectangle (5.880,-1.240);\node at (5.670,-1.470) {0};
    \draw[black,line width=.4pt] (5.880,-1.700) rectangle (6.300,-1.240);\node at (6.090,-1.470) {0};
    \draw[black,line width=.4pt] (6.300,-1.700) rectangle (6.720,-1.240);\node at (6.510,-1.470) {0};
    \draw[line width=.9pt] (3.36,-1.700)--(3.36,-1.240);
  \end{tikzpicture}
  \caption{A byte in memory can contain both committed and unused bits. Dashed cells are not yet part of the stream.}
  \label{fig:revision-h263-bit-window}
\end{figure}

\begin{figure}[!ht]
  \centering
  \begin{tikzpicture}[alt={Native bytes after five guard bytes are BA 35. A solid outline selects those two bytes; all other bytes remain A5.}]
    \node[anchor=north west,inner sep=0,font=\small\bfseries] (title) at (0,0) {Append a field across a byte boundary};
    \node[anchor=north west,inner sep=0] (a) at ([yshift=-2mm]title.south west) {\MARSRevisionCrop{revision/captures/h263-stream/h263-stream-writer.png}{0.99\linewidth}{0.03125}{0.33888889}{0.58203125}{0.60888889}};
    \begin{scope}[shift={(a.south west)},x={($(a.south east)-(a.south west)$)},y={($(a.north west)-(a.south west)$)}]
      \draw[black,line width=.7pt] (0.439063,0.005000) rectangle (0.534687,0.443809);
    \end{scope}
  \end{tikzpicture}
  \caption{The outlined bytes become BA 35 after the nine-bit append at cursor three. Five preceding guard bytes remain A5. The committed count is twelve; the low four bits of the second selected byte remain outside the field.}
  \label{fig:revision-h263-stream-writer}
\end{figure}


The writer's context contains buffer pointer, byte capacity and committed bits.
After validating full-width metadata and $p\leq8c$, compare the requested width
$n$ with $8c-p$. A field that does not fit must fail before buffer access.
The supported width is zero through 32, with no set value bits above that width.
A valid zero-width append needs value zero, succeeds without buffer access and
does not rewrite the context.

For a nonempty admitted field, mask only the selected destination bits. In this
worked fragment, \texttt{t6} initially marks those bits and \texttt{t8} holds
new bits already shifted into position:
\Needspace{10\baselineskip}
\MARSCodeLabel{Replace selected bits while retaining their neighbours}
\begin{minted}{mars-idea}
    nor $t6, $t6, $zero
    lbu $t9, 0($t1)
    and $t9, $t9, $t6
    or $t9, $t9, $t8
    sb $t9, 0($t1)
\end{minted}
The complete writer updates the committed count after output is complete.
A successful field append does not make a sequence of field appends atomic.
A larger structure needs its own plan or staging boundary.

\section{Describe nonzero levels in scan order}

Coefficient arrays use natural row-major order, but the coder visits their
positions in zigzag order. The first natural indices are 0, 1, 8, 16, 9 and 2.
Multiply an index by two to address its halfword. Scan order is a permutation;
it neither changes the transform nor requantizes its values.

For INTRA, emit the separate eight-bit DC code, then scan ordinary AC levels
from position one. For INTER, scan from position zero: its DC coefficient is
an ordinary signed level. Each nonzero event carries RUN, the number of zeros
since the preceding nonzero or scan origin, and LAST, whether any nonzero
follows. Trailing zeros need no events. LAST does not mean that the current
position is the final array element.

Consider INTRA DC code 255, level one at natural index one and minus one at
natural index eight. Both ordinary runs are zero. Their compact descriptions are:
\begin{table}[!ht]
\centering
\begin{tabular}{l r r r l}
\toprule
Part & LAST & RUN & LEVEL & Bits\\
\midrule
DC code & & & 255 & \texttt{11111111}\\
First AC & 0 & 0 & 1 & \texttt{100}\\
Final AC & 1 & 0 & $-1$ & \texttt{01111}\\
\bottomrule
\end{tabular}
\caption{The complete worked block payload occupies sixteen bits.}
\end{table}

The baseline table supplies compact codes for 102 event combinations of LAST,
RUN and magnitude, with a final sign bit. A legal escape instead carries
seven prefix bits \texttt{0000011}, one LAST bit, six RUN bits and an eight-bit
signed LEVEL. Its total is 22 bits. The ordinary escape value cannot be zero or
minus 128. A negative escape level uses two's-complement bytes, rather than the
compact code's separate sign bit.

The same worked block can use its DC byte and two escapes: 52 bits instead of
sixteen. Both representations recover the same levels. An empty INTER block
emits no coefficient payload; an INTRA block still needs its DC byte even when
all AC levels are zero. The macroblock header announces ordinary payload
presence, so an empty block does not need an invented end marker.

Validate all levels and plan the chosen complete payload before touching caller
output. Then stage its touched bytes, preserving the original prefix and unused
suffix bits. Abandon that private work if a child fails. Publish bytes and count
only after the complete block succeeds. Per-field capacity checks alone would
leave a partial block if its last event did not fit.

\begin{figure}[!ht]
  \centering
  \begin{tikzpicture}[alt={After five guard bytes the selected stream bytes are BF F1 E5. The three-byte window has an outline and retains bits outside the sixteen-bit block append.}]
    \node[anchor=north west,inner sep=0,font=\small\bfseries] (title) at (0,0) {Append a complete coefficient block};
    \node[anchor=north west,inner sep=0] (a) at ([yshift=-2mm]title.south west) {\MARSRevisionCrop{revision/captures/h263-stream/h263-stream-coefficients.png}{0.99\linewidth}{0.03125}{0.33888889}{0.58203125}{0.60888889}};
    \begin{scope}[shift={(a.south west)},x={($(a.south east)-(a.south west)$)},y={($(a.north west)-(a.south west)$)}]
      \draw[black,line width=.7pt] (0.439063,0.005000) rectangle (0.582500,0.443809);
    \end{scope}
  \end{tikzpicture}
  \caption{The compact block appends sixteen bits at cursor three: DC code FF, then codes 100 and 01111. Its touched bytes are BF F1 E5. The first three prefix bits and last five unused bits survive; preceding A5 bytes are guards.}
  \label{fig:revision-h263-stream-coefficients}
\end{figure}


\section{Make the macroblock header agree with its payloads}

The decoder needs to know which of the six block payloads follow. Set each
presence bit when that block has an ordinary nonzero level. For INTRA, exclude
the special DC code. Six blocks containing only their required DC values
therefore have ordinary presence pattern zero, although six DC bytes follow.

The four Y blocks form CBPY; the two chroma blocks form CBPC. MCBPC combines
the macroblock type with CBPC. Choose its table group using both picture type
and macroblock type. An INTRA macroblock inside a P-picture does not use the
I-picture INTRA group. The CBPY table uses the direct four-bit pattern for
INTRA and its complement for INTER. A zero INTER Y pattern consequently indexes
fifteen, not zero.

A P-picture macroblock begins with COD. One skips further syntax and means
actual zero motion with no residual levels. Zero introduces the coded header.
A coded INTER macroblock includes MCBPC, CBPY, horizontal and vertical motion
differences, then the present coefficient payloads. An INTRA macroblock has no
motion fields. Its DC bytes remain mandatory regardless of ordinary presence.

Skipping does not mean reuse the predicted motion vector. If the neighbour
predictor is nonzero, coded zero motion and zero motion difference are different
requests. Nor does an unchanged raw block prove a skip: its reconstructed
reference may still differ because of quantization. The reconstruction layer
chooses meaningful levels and vectors; the syntax layer must describe them
without silently substituting another prediction.

\section{Code motion relative to the decoder's predictor}

The search returns full luma pixels. Multiply by two before this syntax boundary,
which uses half-sample components in $[-32,31]$. For selected component $v$ and
predictor $q$, represent their difference modulo 64:
\[
 d=((v-q+32)\bmod64)-32.
\]
With $v=-32$ and $q=31$, the raw difference minus 63 becomes coded difference
one. The decoder adds one to 31 and wraps back to minus 32. A zero coded
difference means selected vector equals predictor, not necessarily zero motion.

The worked table stores one code/width pair per possible difference. Each pair
occupies sixteen bytes, and index $d+32$ can be computed directly:
\Needspace{10\baselineskip}
\MARSCodeLabel{Select a wrapped motion-difference table entry}
\begin{minted}{mars-idea}
    dsubu $t0, $s4, $t0
    daddiu $t0, $t0, 32
    andi $t0, $t0, 63
    dsll $t0, $t0, 4
    dla $t1, syntax_mvd_codes
\end{minted}
Here \texttt{s4} holds the selected component and \texttt{t0} initially holds
its predictor. The masked index is already offset into the table; it is not
the signed difference itself. Complete metadata admission bounds both values
before this calculation.

The general predictor is the component-wise median of left, above and above-right
neighbours, with picture/group border rules. Our controller emits nonempty group
headers between rows. At those row starts the upper candidates are replaced
according to the border rule, so a row begins with predictor zero and subsequent
macroblocks use their preceding actual vector. A skip contributes actual zero.
Do not retain an earlier nonzero predictor as the skipped block's vector.

This illustrates why boundaries belong in the design before optimization.
Group headers change both syntax and neighbour availability. Removing a header
without changing the predictor calculation can make the decoder use another
vector even though the emitted difference code is individually legal.

\section{Publish a whole picture before accepting it as reference}

QCIF has nine rows of eleven macroblocks. For row $r$, column $c$, the Y origin
is $(16c,16r)$ and the chroma origin $(8c,8r)$. Its 594 eight-by-eight blocks
cover $99\cdot6\cdot64=38016$ samples. Copy each reconstructed block row to its
plane stride; a list of block records is not the planar picture layout.

An I-picture has no previous reference. For each later P-picture, this controller
searches the accepted reconstruction, predicts six blocks and encodes their
residuals. In QCIF, a group of blocks, or GOB, contains one macroblock row.
The picture header establishes the first group. Later rows receive
17-bit group start code one, five-bit row number, two-bit group identity and
five-bit quantizer. Reach a byte boundary before each start code using fewer
than eight zero stuffing bits. Do not align every ordinary field or coefficient.

Group identity is constant inside a picture. With the chosen initial I-picture
and later P-pictures, it can be zero for the first and one for every later picture.
It tracks the relevant picture-type change, not alternating picture parity.
The picture header's 22-bit start code has numerical value 32; its temporal
reference has eight bits. The restricted header is fifty bits before macroblocks,
so its following data does not automatically start at the next byte.

\begin{table}[!ht]
\centering
\begin{tabular}{l r l}
\toprule
Header field & Width & Selected profile value\\
\midrule
Picture start code & 22 & 32\\
Temporal reference & 8 & Picture index times six, modulo 256\\
Picture type & 13 & 4160 for I; 4176 for P\\
Quantizer & 5 & 1 through 31; ordinary example eight\\
Multipoint flag & 1 & Zero\\
Extra insertion flag & 1 & Zero\\
\bottomrule
\end{tabular}
\caption{The restricted fifty-bit picture header precedes macroblock data.}
\end{table}

Stage encoded bytes and a complete new reconstruction together. The workspace
reserves 131072 bytes for one encoded picture and 38016 for its reconstructed
planes. The stated worst-case payload and header bound fits that staging arena.
After all children and final alignment succeed, admit caller capacity, copy both
outputs and commit the used count. A late child or capacity failure retains both
caller outputs. Workspace may have changed; it is not a published reference.

There is no need to clear the entire encoded arena before every picture.
Each published bit comes from an emitted field or explicit alignment. Bytes
outside the committed count are not part of the output. This removes repeated
initialization while keeping the serialized meaning exact.

\section{Accept file progress only after both writes complete}

The clip driver owns an input descriptor, a stream descriptor and a local
reconstruction-trace descriptor. Reuse one raw picture buffer, two reconstruction
buffers and the picture workspace. Working storage stays bounded as picture
count grows. After a picture's stream and trace writes both succeed, exchange
old-reference and new-destination roles. Before that point, retain the previous
accepted reference.

\begin{figure}[!ht]
  \centering
  \begin{tikzpicture}[every node/.style={draw,align=center,font=\small,text width=27mm,minimum height=11mm},alt={A staged picture goes through stream write, trace write, then acceptance of the new reference. Failure at either write keeps the preceding accepted reference and reports the accepted prefix, although files can already contain an unaccepted suffix.}]
    \node (a) at (0,0) {Stage picture};
    \node (b) at (3.5,0) {Write stream};
    \node (c) at (7,0) {Write trace};
    \node (d) at (10.5,0) {Accept reference};
    \draw[-{Latex}] (a)--(b);\draw[-{Latex}] (b)--(c);\draw[-{Latex}] (c)--(d);
    \node[text width=75mm] (f) at (5.25,-1.8) {Failure: keep the preceding reference and report the accepted prefix};
    \draw[-{Latex}] (b.south)--(f.north west);\draw[-{Latex}] (c.south)--(f.north east);
  \end{tikzpicture}
  \caption{File writes precede reference promotion. Physical output may include a suffix outside the accepted prefix.}
  \label{fig:revision-h263-accepted-progress}
\end{figure}


Reads and writes return accepted byte counts. Transfer at most 4096 bytes per
request, advance by a positive returned count and retry the suffix. A short
positive result is progress, not completion. Premature read zero, read error or
nonpositive write ends the operation. Close every descriptor actually acquired
on every return path; an earlier successful open still needs cleanup when a
later open fails.

The current educational write service prefetches one extra byte. The worked
transport helper stages a chunk in 4097 private bytes and supplies its guard
there. Only the requested bytes are file content. This isolates the service's
read requirement from exact picture-buffer extents. Changing the request to
include the guard would add unwanted data to the stream.

The progress record contains accepted pictures, accepted encoded bytes and the
last accepted reconstruction pointer. A failed write can leave more bytes in the
file than this record counts. For example, if the second stream write accepts
ten bytes then fails, those bytes are a partial unaccepted suffix. If its stream
write completes but its trace write fails, the second picture is still unaccepted
by this driver. File contents cannot generally be rolled back; progress identifies
the last coherent accepted state.

After the requested pictures, read one additional input byte to require exact EOF.
Trailing data or an error can make the whole operation fail even when the accepted
picture count equals the request. Inspect overall status as well as count.
Metadata rejection precedes path access; path validation precedes opening files.
Compare path contents, not just addresses, so two separate strings cannot name
an input that would then be truncated as an output.

\begin{figure}[!ht]
\centering
\begin{tikzpicture}[alt={Native progress fields show eight accepted pictures, 2212 accepted encoded bytes and a retained reference pointer. The next pair contains copied success zero and accepted count eight. Picture count is outlined and encoded byte count has a dashed outline.}]
\node[anchor=north west,inner sep=0,font=\small\bfseries] (title) at (0,0) {Inspect accepted progress after the complete file writes};
\node[anchor=north west,inner sep=0] (a) at ([yshift=-2mm]title.south west) {\MARSRevisionCrop{revision/captures/h263-stream/h263-stream-clip.png}{0.99\linewidth}{0.03125}{0.29222222}{0.58203125}{0.60888889}};
\begin{scope}[shift={(a.south west)},x={($(a.south east)-(a.south west)$)},y={($(a.north west)-(a.south west)$)}]
\draw[black,line width=.7pt] (0.582500,0.476910) rectangle (0.965000,0.704865);
\draw[black,dashed,line width=.7pt] (0.200000,0.240955) rectangle (0.582500,0.468910);
\end{scope}
\end{tikzpicture}
\caption{After a leading reserved zero, the three progress fields are accepted count eight, encoded bytes 2212 (dashed), and the last reference pointer. The final row contains the caller copies of success zero and count eight. Both required files have been written and all owned descriptors closed.}
\label{fig:revision-h263-stream-clip}
\end{figure}


\section{Check the actual stream with independent tools}

From the book repository, prepare and run the completed source in a new directory:
\Needspace{10\baselineskip}
\MARSCodeLabel{Generate input, execute the MIPS encoder and check its files}
\begin{Verbatim}
python3 packages/035-h263/construction/h263-clip-construction/run.py \
  revision/labs/files/h263-encoder.asm --output-dir work/run-eight \
  --cli "$CLI" --java "$JAVA"
\end{Verbatim}
The helper generates raw pixels and installs paths/controls while retaining the
instruction section. MARS produces \texttt{encoded.263} and
\texttt{reconstructed.yuv}. FFmpeg is used afterward as a decoder, not as a
host encoder supplying the student's answer. The directory also retains raw
input, metadata, prepared source, execution log and verification report.

The strict parser visits every picture, group, macroblock and coefficient payload.
It checks the restricted profile and temporal references, rather than recognizing
only a plausible first header. FFmpeg then decodes every picture. Compare decoded
samples with raw input to measure quality and with local reconstruction to measure
decoder agreement. These are different questions. The assignment requires at least
35 dB PSNR for each non-identical picture plane and at most sixteen sample values
of local/decoder difference, alongside the separate kernel and inverse-accuracy
gates. Those clip limits are not substitutes for accurate inverse arithmetic.

PSNR is peak signal-to-noise ratio. For $n$ samples, first compute mean squared
error $\operatorname{MSE}=\sum_i(S_i-X_i)^2/n$. This check uses the full eight-bit
peak value 255:
\[
 \operatorname{PSNR}=10\log_{10}\!\left(\frac{255^2}{\operatorname{MSE}}\right).
\]
Higher PSNR means less mean squared error. An exactly matching plane has zero
MSE and unbounded PSNR; the report treats that as an exact match. Evaluate planes
separately so a large luma plane cannot conceal a poor chroma result. Average
error can also hide an isolated large difference. The maximum local/decoder
difference examines agreement between those two reconstructions, rather than
quality relative to the raw source.

The renewed eight-picture reference produces 2212 bytes; the equivalent
escape-oriented implementation produces 5101. Both independently decode with
maximum local/decoder difference two. The 44-picture reference run also crosses
temporal-reference wraparound and remains within the input limit. These are
observed files, not required private bit strings for every valid encoder.

Keep picture time separate from execution time. Temporal reference advances by
six ticks modulo 256 under a $30000/1001$-tick clock. Eight supplied intervals
last 1.6016 seconds; this includes the fixture's final picture hold, not just its
last start time. Forty-four pictures end with reference two after crossing 255.
The host can take longer than the represented video to encode it.

The ordinary reference MSA run executes 85093476 guest instructions, compared
with 206396173 in scalar mode. Those complete-run observations include assembly
expansions, file helpers and caller work; exact counts also depend on prepared
paths. A smaller stream does not necessarily require fewer instructions, and an
uncontrolled whole-JVM wall time is not a simulator throughput benchmark. Retain
source, input, tool identities and the measurement boundary when comparing changes.

\section{Read a bounded field without losing the cursor}

Implement \texttt{h263\_read\_bits} in the independent lab file. Its context
separates byte capacity, valid bit count and read cursor. The unused low bits of
\texttt{BA 35} remain present in memory, but with twelve valid bits they are not
available as another field. Starting at cursor three, reading nine bits returns
419 and advances to twelve. A further one-bit request must fail without advancing.
A valid empty read returns zero and leaves the cursor untouched.

For one bit at position $p$, load byte $\lfloor p/8\rfloor$, select
$7-(p\bmod8)$, shift the result accumulator left and append that bit. A second
implementation can consume the available fragment of each byte. Both must produce
a zero-extended result up to 32 bits, preserve input and publish the cursor only
after the entire requested field is admitted. A failed or empty call must not
reuse a previous nonzero return value.

The lab gives complete metadata, access and failure contracts, plus checks that
connect the reader to actual MARS-written fields. Once the primitive works, use
it to read the first 22-bit picture start code and eight-bit temporal reference.
Do not call that small extension a complete decoder: later variable-length events,
prediction and reconstructed samples have their own meanings and boundaries.

\begin{figure}[!ht]
\centering
\begin{tikzpicture}[alt={The caller observation pairs are success zero and returned field 419; cursor twelve and valid bit count twelve. Field 419 is outlined and cursor twelve has a dashed outline.}]
\node[anchor=north west,inner sep=0,font=\small\bfseries] (title) at (0,0) {Separate returned value from committed read cursor};
\node[anchor=north west,inner sep=0] (a) at ([yshift=-2mm]title.south west) {\MARSRevisionCrop{revision/captures/h263-stream/h263-stream-reader.png}{0.99\linewidth}{0.03125}{0.31555556}{0.58203125}{0.60888889}};
\begin{scope}[shift={(a.south west)},x={($(a.south east)-(a.south west)$)},y={($(a.north west)-(a.south west)$)}]
\draw[black,line width=.7pt] (0.582500,0.313824) rectangle (0.965000,0.614647);
\draw[black,dashed,line width=.7pt] (0.200000,0.005000) rectangle (0.582500,0.305824);
\end{scope}
\end{tikzpicture}
\caption{The completed demo returns 419 (outlined) and advances its cursor to twelve (dashed). The final field copies valid count twelve. The caller obtains these values from the actual function and context; unused low bits of BA 35 are unavailable to a later field.}
\label{fig:revision-h263-stream-reader}
\end{figure}


\section*{Sources and continuing work}

The syntax and prediction rules follow clauses 4.5, 5 and 6.1 of
\href{https://www.itu.int/rec/T-REC-H.263-199802-S/en}{ITU-T H.263 (02/1998)}.
The public profile and callable contracts document our narrower encoder policy.
The worked file and construction package provide a starting point for other
legal coding choices. Keep independent decoding, quality, inverse accuracy,
retained-state behavior and measured work together when developing them.
